%! Linear Transformations — Unit 5, Lesson 5.3 — Student Packet Cover Sheet
\documentclass[10pt]{article}
\usepackage{linalg-article}
\usepackage{linalg-boxes}
\usepackage{ltablex}
\keepXColumns

\begin{document}

% Full-bleed burgundy banner
\begin{tikzpicture}[remember picture, overlay]
  \fill[burgundy] (current page.north west)
    rectangle ([yshift=-0.9in]current page.north east);
\end{tikzpicture}
\vspace{-0.8in}

\begin{center}
  {\color{white}\textbf{\LARGE Linear Algebra}}\\[4pt]
  {\color{blushmid}\large\textbf{Unit 5 \quad Determinants and Linear Transformations}}\\[6pt]
  {\color{white}\textbf{\large Lesson 5.3 \quad Linear Transformations}}
\end{center}

\vspace{0.22in}
\namedateperiod
\vspace{0.18in}

\begin{learningtargetbox}
\small
\begin{enumerate}[leftmargin=1.5em, itemsep=5pt]
  \item \ldots{} read a matrix as a \textbf{motion of the plane} $\mathbf{x}\mapsto A\mathbf{x}$, and find where the unit square lands (the parallelogram of the columns).
  \item \ldots{} name the motion --- \textbf{stretch, rotation, reflection, shear} --- from the matrix.
  \item \ldots{} use $|\det A|$ as the factor \textbf{every area is multiplied by}, with the \textbf{sign} of $\det A$ reporting orientation and $\det A=0$ meaning the map \textbf{collapses} (no inverse).
\end{enumerate}
\end{learningtargetbox}

\vspace{0.14in}

\begin{tocbox}
\small
\renewcommand{\arraystretch}{1.5}
\begin{tabularx}{\linewidth}{@{} c l X r @{}}
\rowcolor{blush}
\textbf{\#} & \textbf{Component} & \textbf{Description} & \textbf{Score} \\
\midrule
1 & Warm-Up        & Spiral: $A\mathbf{x}$ as a column combination, a $2\times2$ determinant, and area & \blank{1.2cm} \\
2 & Guided Notes   & Matrices as motions; unit square $\to$ column parallelogram; $|\det A|$ as area factor & \blank{1.2cm} \\
3 & Group Activity & Moving Space with Matrices --- naming motions, area factors, orientation             & \blank{1.2cm} \\
4 & Exit Ticket    & Quick check on area factor, orientation, and a $\det=0$ collapse                     & \blank{1.2cm} \\
5 & Homework       & Image areas, motions and orientation, composition, and a $3$-D volume factor        & \blank{1.2cm} \\
\midrule
  & \multicolumn{2}{r}{\textbf{Total}} & \blank{1.2cm} \\
\end{tabularx}
\end{tocbox}

\vspace{0.10in}

\begin{remindbox}
\small A matrix does not just store numbers --- it \textbf{moves space}. Multiplying by $A$ sends the
unit square to the parallelogram spanned by the \emph{columns} of $A$, and keeps grid lines straight, so
\emph{every} shape is transformed the same way. The area of that image parallelogram is exactly
$|\det A|$, so applying $A$ multiplies every area by $|\det A|$. The \textbf{sign} of $\det A$ tells you
whether the picture was flipped (orientation), and $\det A=0$ means space was \textbf{collapsed} flat ---
which is why such a matrix has no inverse.
\end{remindbox}

\end{document}
